% =====================================================================
\chapter{Sumatorios, productorios y binomio de Newton}
\label{tema:4}
% =====================================================================

La probabilidad de una secuencia es un producto de probabilidades condicionadas, su logaritmo es una suma, y la pérdida de entrenamiento es una media sobre millones de tokens. Manejar con soltura la notación $\sum$ y $\prod$ es, por tanto, imprescindible. Este tema sigue el capítulo~2 «Sums» de Graham, Knuth y Patashnik~\cite{gkp} y el capítulo~14 «Sums and Asymptotics» de Lehman, Leighton y Meyer~\cite{mcs}; el binomio de Newton sigue el §9.6 de OpenStax~\cite{openstax}.

\section{Notación y principio de inducción}

\begin{definicion}[Sumatorio y productorio]
Dados enteros $m\le n$ y números $a_m,\dots,a_n$,
\[
\sum_{k=m}^{n}a_k=a_m+a_{m+1}+\dots+a_n,\qquad \prod_{k=m}^{n}a_k=a_m\,a_{m+1}\cdots a_n.
\]
La letra $k$ es un \emph{índice mudo}: puede sustituirse por cualquier otra sin cambiar el valor. Si $n<m$, la suma es \emph{vacía} y vale~0, y el producto vacío vale~1. Más en general, $\sum_{k\in S}a_k$ denota la suma sobre un conjunto finito $S$ de índices. El \emph{factorial} es $n!=\prod_{k=1}^{n}k$, con $0!=1$.
\end{definicion}

Formalmente, la suma se define por recurrencia: $\sum_{k=m}^{m-1}a_k=0$ y $\sum_{k=m}^{n}a_k=\sum_{k=m}^{n-1}a_k+a_n$. Por eso las fórmulas para sumas se demuestran de forma natural por inducción.

\begin{teorema}[Principio de inducción]\label{teo:induccion}
Sea $P(n)$ una afirmación sobre los enteros $n\ge n_0$. Si $P(n_0)$ es cierta y, para cada $n\ge n_0$, $P(n)$ implica $P(n+1)$, entonces $P(n)$ es cierta para todo $n\ge n_0$.
\end{teorema}

Este principio es equivalente al principio del buen orden (todo subconjunto no vacío de $\N$ tiene mínimo); ambos se estudian en~\cite[cap.~2 y~5]{mcs} y~\cite[cap.~10]{hammack}.

\section{Propiedades de los sumatorios}

\begin{proposicion}\label{prop:sumas}
Para números $a_k,b_k$ y una constante $c$:
\begin{enumerate}
  \item linealidad: $\sum_{k=m}^n(ca_k+b_k)=c\sum_{k=m}^na_k+\sum_{k=m}^nb_k$;
  \item partición: si $m\le p<n$, $\sum_{k=m}^na_k=\sum_{k=m}^pa_k+\sum_{k=p+1}^na_k$;
  \item suma de una constante: $\sum_{k=m}^nc=(n-m+1)\,c$;
  \item cambio de índice: $\sum_{k=m}^na_k=\sum_{j=m+r}^{n+r}a_{j-r}$ para todo entero $r$, y $\sum_{k=0}^na_k=\sum_{k=0}^na_{n-k}$.
\end{enumerate}
\end{proposicion}

\begin{proof}
(i)--(iii) son consecuencia de la asociatividad, la conmutatividad y la distributividad de los axiomas de cuerpo, aplicadas un número finito de veces (formalmente, por inducción sobre el número de términos). (iv) En ambos casos se suman exactamente los mismos números en distinto orden: con $j=k+r$, cuando $k$ recorre $m,\dots,n$, $j$ recorre $m+r,\dots,n+r$; con $j=n-k$, cuando $k$ recorre $0,\dots,n$, $j$ recorre $n,\dots,0$.
\end{proof}

No hay reglas análogas para el producto: en general $\sum a_kb_k\neq\bigl(\sum a_k\bigr)\bigl(\sum b_k\bigr)$.

\section{Sumas clásicas}

\begin{teorema}[Suma aritmética]\label{teo:aritmetica}
Para todo $n\in\N$, $\displaystyle\sum_{k=1}^nk=\frac{n(n+1)}2$. En general, si $a_k=a_1+(k-1)d$, entonces $\sum_{k=1}^na_k=\dfrac{n(a_1+a_n)}{2}$.
\end{teorema}

\begin{proof}
Sea $S=\sum_{k=1}^nk$. Por la \cref{prop:sumas}(iv), $S=\sum_{k=1}^n(n+1-k)$. Sumando ambas expresiones, $2S=\sum_{k=1}^n(n+1)=n(n+1)$. La fórmula general se obtiene igual, porque $a_k+a_{n+1-k}=a_1+a_n$ para todo $k$.
\end{proof}

\begin{teorema}[Suma geométrica]\label{teo:geometrica}
Para $r\neq1$ y $n\in\N_0$, $\displaystyle\sum_{k=0}^nr^k=\frac{1-r^{n+1}}{1-r}$. Si $r=1$, la suma vale $n+1$.
\end{teorema}

\begin{proof}
Por el \emph{método de perturbación}~\cite[§2.3]{gkp}: sea $S=\sum_{k=0}^nr^k$. Entonces $rS=\sum_{k=1}^{n+1}r^k$ y, restando, se cancelan todos los términos salvo dos: $S-rS=1-r^{n+1}$. Como $r\neq1$, se divide por $1-r$.
\end{proof}

\begin{proposicion}[Sumas telescópicas]\label{prop:telescopica}
Para cualesquiera números $a_m,\dots,a_{n+1}$,
\[
\sum_{k=m}^n(a_{k+1}-a_k)=a_{n+1}-a_m.
\]
\end{proposicion}

\begin{proof}
Por linealidad y cambio de índice, $\sum_{k=m}^na_{k+1}-\sum_{k=m}^na_k=\sum_{j=m+1}^{n+1}a_j-\sum_{k=m}^na_k=a_{n+1}-a_m$, ya que los términos con índice entre $m+1$ y $n$ aparecen en ambas sumas.
\end{proof}

\begin{ejemplo}\label{ej:telescopica}
Como $\frac1{k(k+1)}=\frac1k-\frac1{k+1}$, se tiene $\displaystyle\sum_{k=1}^n\frac1{k(k+1)}=1-\frac1{n+1}$.
\end{ejemplo}

\begin{teorema}[Sumas de potencias]\label{teo:potencias}
Para todo $n\in\N$,
\[
\sum_{k=1}^nk^2=\frac{n(n+1)(2n+1)}6,\qquad\sum_{k=1}^nk^3=\Bigl(\frac{n(n+1)}2\Bigr)^2.
\]
\end{teorema}

\begin{proof}
Para la primera, se suma la identidad $(k+1)^3-k^3=3k^2+3k+1$ para $k=1,\dots,n$. El miembro izquierdo es telescópico y vale $(n+1)^3-1$; el derecho vale $3S_2+3\frac{n(n+1)}2+n$, donde $S_2=\sum k^2$. Despejando,
\[
3S_2=(n+1)^3-1-\tfrac32n(n+1)-n=\tfrac{(n+1)}{2}\bigl[2(n+1)^2-3n-2\bigr]=\tfrac{(n+1)n(2n+1)}{2}.
\]
La segunda fórmula se demuestra por inducción en el \cref{eje:4-cubos}.
\end{proof}

\begin{ejemplo}[Medias exponenciales]\label{ej:ema}
El optimizador Adam (módulo~7) mantiene la media móvil $m_t=\beta m_{t-1}+(1-\beta)g_t$ con $m_0=0$ y $0<\beta<1$. Desplegando la recurrencia, $m_t=(1-\beta)\sum_{k=1}^{t}\beta^{t-k}g_k$. Si todos los gradientes valen $g$, el \cref{teo:geometrica} da $m_t=(1-\beta)g\sum_{j=0}^{t-1}\beta^{j}=(1-\beta^t)\,g$: la media está sesgada hacia cero. Con $\beta=0{,}9$ y $t=10$, $m_{10}\approx0{,}651\,g$. Adam divide por $1-\beta^t$ para corregir ese sesgo.
\end{ejemplo}

\section{Sumas dobles}

Una suma doble $\sum_{i}\sum_{j}a_{ij}$ suma los elementos de una tabla. Si los índices recorren un rectángulo, el orden de suma puede intercambiarse:
\[
\sum_{i=1}^m\sum_{j=1}^na_{ij}=\sum_{j=1}^n\sum_{i=1}^ma_{ij}.
\]
Si recorren un triángulo, hay que reescribir los límites~\cite[§2.4]{gkp}:
\[
\sum_{i=1}^n\sum_{j=i}^na_{ij}=\sum_{1\le i\le j\le n}a_{ij}=\sum_{j=1}^n\sum_{i=1}^ja_{ij}.
\]

\begin{proposicion}\label{prop:producto-sumas}
$\displaystyle\Bigl(\sum_{i=1}^ma_i\Bigr)\Bigl(\sum_{j=1}^nb_j\Bigr)=\sum_{i=1}^m\sum_{j=1}^na_ib_j$. En particular, $\displaystyle\Bigl(\sum_{i=1}^na_i\Bigr)^2=\sum_{i=1}^na_i^2+2\sum_{1\le i<j\le n}a_ia_j$.
\end{proposicion}

\begin{proof}
Por linealidad, $\bigl(\sum_ia_i\bigr)B=\sum_ia_iB$ con $B=\sum_jb_j$, y de nuevo $a_iB=\sum_ja_ib_j$. Para la segunda, con $b=a$ se separan en la suma doble los términos con $i=j$, los de $i<j$ y los de $i>j$; los dos últimos grupos son iguales por simetría.
\end{proof}

\begin{ejemplo}[Coste de la atención causal]\label{ej:causal}
En un Transformer con máscara causal (módulo~13), el token $i$ atiende a los tokens $j\le i$. El número de pares $(i,j)$ con $1\le j\le i\le n$ es $\sum_{i=1}^n\sum_{j=1}^i1=\sum_{i=1}^ni=\frac{n(n+1)}2$, aproximadamente la mitad de los $n^2$ pares de la atención completa.
\end{ejemplo}

\section{Productorios y binomio de Newton}

Los productorios cumplen $\prod_k(a_kb_k)=\prod_ka_k\prod_kb_k$, $\prod_{k=1}^nc=c^n$ y, si todos los $a_k$ son positivos, $\ln\prod_ka_k=\sum_k\ln a_k$ (\cref{tema:5}).

\begin{definicion}[Números combinatorios]
Para enteros $0\le k\le n$, $\displaystyle\binom nk=\frac{n!}{k!\,(n-k)!}$, y $\binom nk=0$ si $k<0$ o $k>n$.
\end{definicion}

\begin{lema}[Identidad de Pascal]\label{lem:pascal}
Para $1\le k\le n$, $\displaystyle\binom nk+\binom n{k-1}=\binom{n+1}k$.
\end{lema}

\begin{proof}
Con denominador común $k!\,(n+1-k)!$:
\[
\frac{n!}{k!(n-k)!}+\frac{n!}{(k-1)!(n-k+1)!}=\frac{n!\,\bigl[(n+1-k)+k\bigr]}{k!\,(n+1-k)!}=\frac{(n+1)!}{k!\,(n+1-k)!}.\qedhere
\]
\end{proof}

\begin{teorema}[Binomio de Newton]\label{teo:binomio}
Para todos $a,b\in\R$ y $n\in\N_0$, $\displaystyle(a+b)^n=\sum_{k=0}^n\binom nka^{n-k}b^k$.
\end{teorema}

\begin{proof}
Por inducción sobre $n$. Para $n=0$ ambos miembros valen~1. Supóngase cierto para $n$. Entonces
\[
(a+b)^{n+1}=(a+b)\sum_{k=0}^n\binom nka^{n-k}b^k=\sum_{k=0}^n\binom nka^{n+1-k}b^k+\sum_{k=1}^{n+1}\binom n{k-1}a^{n+1-k}b^k,
\]
donde en la segunda suma se ha cambiado $k$ por $k-1$. Usando el convenio $\binom n{-1}=\binom n{n+1}=0$, ambas sumas pueden extenderse a $k=0,\dots,n+1$, y por el \cref{lem:pascal} el coeficiente de $a^{n+1-k}b^k$ es $\binom nk+\binom n{k-1}=\binom{n+1}k$ (para $k=0$ y $k=n+1$ vale~1, que también coincide).
\end{proof}

\begin{corolario}\label{cor:binomio}
$\displaystyle\sum_{k=0}^n\binom nk=2^n$ y, si $n\ge1$, $\displaystyle\sum_{k=0}^n(-1)^k\binom nk=0$.
\end{corolario}

\begin{proof}
Basta tomar $a=b=1$ y $a=1$, $b=-1$ en el \cref{teo:binomio}.
\end{proof}

\begin{observacion}[Errores frecuentes]
\begin{enumerate}[label=(\alph*)]
\item $\sum_{k=1}^n c=nc$, no $c$.
\item $\sum a_kb_k\neq\sum a_k\sum b_k$ y $\bigl(\sum a_k\bigr)^2\neq\sum a_k^2$ (\cref{prop:producto-sumas}).
\item Errores de uno en los límites: $\sum_{k=m}^n$ tiene $n-m+1$ términos, no $n-m$.
\item La fórmula geométrica exige $r\neq1$; al cambiar el índice hay que cambiar a la vez los límites y el término general.
\end{enumerate}
\end{observacion}

\begin{observacion}[Conexión con el curso]
Por la regla de la cadena (módulo~1), la probabilidad de una secuencia es $\prod_tP(x_t\mid x_{<t})$ y su logaritmo es $\sum_t\ln P(x_t\mid x_{<t})$: los productorios se convierten en sumatorios. La serie geométrica reaparece en el retorno descontado $\sum_t\gamma^tr_t$ del aprendizaje por refuerzo (módulo~21), y los números combinatorios en el estimador pass@$k$ del módulo~22.
\end{observacion}

\begin{lecturas}
\item Graham, Knuth y Patashnik~\cite{gkp}, §2.1 «Notation», §2.3 «Manipulation of Sums» (método de perturbación) y §2.4 «Multiple Sums».
\item Lehman, Leighton y Meyer~\cite{mcs}, §14.1 «The Value of an Annuity», §14.2 «Sums of Powers», §14.5 «Products» y §14.6 «Double Trouble».
\item OpenStax \emph{College Algebra}~\cite{openstax}, §9.2--9.4 (sucesiones aritméticas y geométricas, series) y §9.6 «Binomial Theorem».
\item Hammack~\cite{hammack}, cap.~10 «Mathematical Induction».
\end{lecturas}

\section*{Ejercicios del tema 4}
\addcontentsline{toc}{section}{Ejercicios del tema 4}

\begin{ejercicio}{1}{eje:4-notacion}
Escriba con notación de sumatorio: (a) $1+3+5+\dots+99$; (b) $2+4+8+\dots+1024$; (c) $\frac12-\frac13+\frac14-\dots+\frac1{100}$.
\end{ejercicio}

\begin{ejercicio}{1}{eje:4-calcular}
Calcule: (a) $\sum_{k=1}^4k^2$; (b) $\prod_{k=1}^5\frac k{k+1}$; (c) $\sum_{k=4}^{0}q_k$ con el convenio de este tema; (d) $\sum_{k=1}^{10}7$.
\fuente{adaptado de Graham, Knuth y Patashnik~\cite{gkp}, cap.~2, ejercicios de calentamiento (convenio sobre sumas vacías).}
\end{ejercicio}

\begin{ejercicio}{1}{eje:4-errores}
Encuentre el error en cada «igualdad»: (a) $\sum_{k=1}^n3=3$; (b) $\sum_{k=1}^2a_kb_k=(a_1+a_2)(b_1+b_2)$; (c) $\bigl(\sum_{k=1}^2a_k\bigr)^2=a_1^2+a_2^2$. Dé un contraejemplo numérico de (b).
\end{ejercicio}

\begin{ejercicio}{2}{eje:4-aritmeticas}
(a) Calcule $\sum_{k=1}^{100}(3k-2)$. (b) Calcule la suma de todos los múltiplos de~7 comprendidos entre 100 y 1000.
\fuente{adaptado de OpenStax \emph{College Algebra}~\cite{openstax}, §9.2 y §9.4.}
\end{ejercicio}

\begin{ejercicio}{2}{eje:4-geometricas}
Calcule: (a) $\sum_{k=0}^{9}3\cdot2^k$; (b) $\sum_{k=1}^{n}\bigl(\frac12\bigr)^k$; (c) $\sum_{k=0}^{n}(-1)^kx^k$ para $x\neq-1$.
\fuente{adaptado de OpenStax \emph{College Algebra}~\cite{openstax}, §9.3--9.4.}
\end{ejercicio}

\begin{ejercicio}{2}{eje:4-telescopicas}
Calcule en forma cerrada: (a) $\sum_{k=1}^n\frac1{k(k+2)}$ y su valor para $n=10$; (b) $\sum_{k=1}^n\ln\bigl(1+\frac1k\bigr)$; (c) $\sum_{k=1}^{99}\frac1{\sqrt k+\sqrt{k+1}}$.
\end{ejercicio}

\begin{ejercicio}{2}{eje:4-indice}
(a) Escriba $\sum_{k=3}^{n+2}(k-2)^2$ como una suma que empiece en~1 y calcúlela. (b) Calcule $\sum_{k=1}^nk(k+1)$ en forma cerrada usando las fórmulas del tema y compruébela para $n=10$.
\end{ejercicio}

\begin{ejercicio}{2}{eje:4-dobles}
Calcule: (a) $\sum_{i=1}^3\sum_{j=1}^4(i+2j)$; (b) $\sum_{i=1}^n\sum_{j=1}^ij$; (c) intercambie el orden de suma en $\sum_{i=1}^n\sum_{j=i}^n\frac1j$ y calcule su valor.
\fuente{adaptado de Lehman, Leighton y Meyer~\cite{mcs}, §14.6 «Double Trouble».}
\end{ejercicio}

\begin{ejercicio}{2}{eje:4-binomio}
(a) Desarrolle $(2x-1)^5$. (b) Halle el coeficiente de $x^4$ en $\bigl(x+\frac2x\bigr)^8$. (c) Calcule $\sum_{k=0}^{6}\binom6k2^k$ sin sumar término a término.
\fuente{adaptado de OpenStax \emph{College Algebra}~\cite{openstax}, §9.6.}
\end{ejercicio}

\begin{ejercicio}{2}{eje:4-productos}
Calcule: (a) $\prod_{k=2}^n\bigl(1-\frac1{k^2}\bigr)$ y su valor para $n=10$; (b) $\prod_{k=1}^n2^k$; (c) $\ln\prod_{k=1}^{n}\frac{k+1}{k}$.
\fuente{adaptado de Lehman, Leighton y Meyer~\cite{mcs}, §14.5 «Products».}
\end{ejercicio}

\begin{ejercicio}{3}{eje:4-cubos}
Demuestre por inducción que $\displaystyle\sum_{k=1}^nk^3=\Bigl(\frac{n(n+1)}2\Bigr)^2$ para todo $n\in\N$.
\fuente{adaptado de Lehman, Leighton y Meyer~\cite{mcs}, problema~2.17 (allí se pide por el principio del buen orden).}
\end{ejercicio}

\begin{ejercicio}{3}{eje:4-perturbacion}
Use el método de perturbación para demostrar que, si $r\neq1$,
\[
\sum_{k=0}^nkr^k=\frac{r-(n+1)r^{n+1}+nr^{n+2}}{(1-r)^2}.
\]
\fuente{adaptado de Graham, Knuth y Patashnik~\cite{gkp}, §2.3.}
\end{ejercicio}

\begin{ejercicio}{3}{eje:4-kbinom}
Demuestre que $k\binom nk=n\binom{n-1}{k-1}$ para $1\le k\le n$ y deduzca que $\sum_{k=0}^nk\binom nk=n\,2^{n-1}$.
\end{ejercicio}

\begin{ejercicio}{3}{eje:4-basilea}
Demuestre por inducción que $\displaystyle\sum_{k=1}^n\frac1{k^2}\le2-\frac1n$ para todo $n\in\N$.
\fuente{adaptado de Hammack~\cite{hammack}, cap.~10, ejercicios.}
\end{ejercicio}

\begin{ejercicio}{3}{eje:4-varianza}
Sean $a_1,\dots,a_n\in\R$. Demuestre que
\[
\sum_{i=1}^n\sum_{j=1}^n(a_i-a_j)^2=2n\sum_{i=1}^na_i^2-2\Bigl(\sum_{i=1}^na_i\Bigr)^2
\]
y deduzca que $\bigl(\sum_ia_i\bigr)^2\le n\sum_ia_i^2$, con igualdad si y solo si todos los $a_i$ son iguales.
\end{ejercicio}
